\chapter{Keyboard parsing and CSI-u}
\label{ch:keyboard}

\begin{aside}
Pressing Ctrl-Shift-K should be distinguishable from Ctrl-K.
Pressing the dedicated copy key on your keyboard should not look
like a stream of random letters. The terminal world has been
working on this since 1984; the answer arrived in 2017.
\end{aside}

\section{The basic key encoding}

Most keys arrive as their literal character: typing \code{a}
sends the byte \code{0x61}. Modifiers fold the byte:

\begin{itemize}[leftmargin=*]
  \item \textbf{Shift:} the shifted character (\code{A}).
  \item \textbf{Ctrl:} the byte minus 64, masked to 5 bits.
    \code{Ctrl-A} sends \code{0x01}, \code{Ctrl-C} sends
    \code{0x03}, etc.
  \item \textbf{Alt:} the character with an \code{ESC} prefix.
    \code{Alt-a} sends \code{ESC a} (two bytes).
\end{itemize}

\section{Function keys and arrows}

These send escape sequences:

\begin{itemize}[leftmargin=*]
  \item Arrow up: \code{ESC [ A} (or \code{ESC O A} in
    application mode).
  \item F1--F4: \code{ESC O P} through \code{ESC O S}.
  \item F5--F12: \code{ESC [ 15~} through \code{ESC [ 24~}.
  \item Home, End, PgUp, PgDn: various \code{ESC [ ... ~}.
\end{itemize}

The variation between terminals here is enormous. \maya{} ships
a table of known sequences (terminfo-style) and falls back to
``unknown escape'' for anything else.

\section{The ambiguity problem}

\code{Esc} alone is a real key. \code{Alt-x} also begins with
\code{ESC}. How does the parser tell them apart?

The convention: if a byte arrives within a short window after
\code{ESC} (typically 50 ms), it is an \code{Alt-}prefixed key.
Otherwise it is a bare \code{Esc}.

\maya{} uses a 50 ms timeout. The cost: pressing \code{Esc} adds
50 ms before it dispatches. Most users do not notice; vim does
the same.

\section{Ctrl-letter collisions}

\code{Ctrl-I} is \code{Tab}. \code{Ctrl-M} is \code{Enter}.
\code{Ctrl-[} is \code{Esc}. \code{Ctrl-H} is \code{Backspace}
on some terminals.

These collisions are not a parser bug --- they are baked into
ASCII. \code{Ctrl-letter} sends the control character, and
control characters have their own purposes.

\maya{} reports the byte and the user's keymap decides what to
do. ``Bind \code{Ctrl-I} to something other than tab'' is not
possible at the byte layer.

\section{CSI-u: the modern answer}

Paul Evans's CSI-u protocol (2017) gives each key a unique
encoding:

\begin{lstlisting}
ESC [ code ; modifiers u
\end{lstlisting}

Where \code{code} is the Unicode code point of the unshifted
key and \code{modifiers} is a bitmask (1 = shift, 2 = alt,
4 = ctrl, 8 = super).

So \code{Ctrl-Shift-K} arrives as \code{ESC [ 107 ; 5 u} ---
unambiguous, full modifier mask, no collisions.

\maya{} requests CSI-u at startup:

\begin{lstlisting}
emit("\x1b[>1u");   // enable CSI-u, level 1
\end{lstlisting}

And on shutdown:

\begin{lstlisting}
emit("\x1b[<u");    // disable
\end{lstlisting}

Terminals that don't support CSI-u ignore the request and keep
sending old-style sequences.

\section{Kitty keyboard protocol}

The Kitty terminal extends CSI-u further: separate codes for
key press and release, repeat events, alternate-key reporting
(distinguish \code{a} from \code{shift-a} even though they
produce the same character), and a more granular modifier set
(super, hyper, meta, caps lock state).

\maya{} supports Kitty's protocol when detected. The detection
is a query (\code{ESC [ ? u}) and a parse of the response.

\section{Modifier-only keys}

Holding \code{Shift} alone with no other key does not send
anything on most terminals. CSI-u level 4 lets it; Kitty's
protocol supports it.

Modal apps that want to show ``shift is held'' (drag-and-drop
indicators, for instance) need a CSI-u level 4 terminal or
must do without.

\section{Paste}

A paste of many characters arrives as fast as the terminal can
write to the pty. Bracketed paste mode wraps it:

\begin{lstlisting}
ESC [ 200 ~
... pasted bytes ...
ESC [ 201 ~
\end{lstlisting}

\maya{} requests bracketed paste at startup (\code{ESC [ ?
2004 h}) and the parser groups the bytes between the brackets
into a single \code{Paste} event. The app can then handle it
as one operation, not as a flood of key events.

\begin{warning}
Without bracketed paste, a pasted shell command can fire
keybinds in your app. With it, the paste is one event and the
keybinds are untouched.
\end{warning}

\section{Composing characters}

On macOS, holding a letter offers accents in a small popup.
The terminal does not see the popup --- it sees the final
character once you commit. Composition is OS-level, not
terminal-level.

For East Asian input methods (IME) the same applies: \maya{}
sees the committed string, not the candidate window.

\section{Key event type}

\maya{}'s normalised key event:

\begin{lstlisting}
struct KeyEvent {
    char32_t code;          // U+0000 for special keys
    enum Special {
        None, Up, Down, Left, Right,
        Home, End, PgUp, PgDn,
        F1, F2, /* ... */ F12,
        Enter, Tab, Escape, Backspace, Delete,
    } special;
    bool shift, alt, ctrl, super;
};
\end{lstlisting}

The dispatcher walks focus order and asks each focused
element ``do you want this key?''.

\section{Recap}

\begin{recap}
\begin{itemize}[leftmargin=*]
  \item Basic keys: literal byte; Ctrl folds, Alt prefixes.
  \item CSI-u gives every key a unique encoding.
  \item Bracketed paste groups pasted bytes into one event.
  \item Modifier-only keys need CSI-u level 4.
  \item Composition (accents, IME) happens before \maya{}
    sees anything.
\end{itemize}
\end{recap}

\section*{Workshop}

\begin{enumerate}[leftmargin=*]
  \item Print every key event to a side panel. Confirm
    CSI-u shows up by pressing \code{Ctrl-Shift-K}.
  \item Add a paste handler that inserts the pasted block
    as a single undo operation.
  \item Detect bare \code{Esc} vs \code{Alt-x} and tune the
    timeout. What feels best?
\end{enumerate}
